\medskip












\medskip

\begin{enumerate}
\item Une primitive de la fonction $f$, définie sur $\R$ par $f(x) =x\text{e}^x$, est la fonction $F$, définie sur $\R$, par :






Avec $F(x) = (x - 1)\text{e}^x$, on a $F'(x) = 1\text{e}^x + (x - 1)\text{e}^x = \text{e}^x(1 + x - 1) = x\text{e}^x$. Réponse \textbf{b.}

\item On considère la fonction $g$ définie par $g(x) = \ln \left(\dfrac{x - 1}{2x+ 4}\right).$









La fonction est définie quand l'argument du logarithme est supérieur à zéro donc quand 

$\dfrac{x - 1}{2x+ 4} > 0$ ou ce qui revient au même quand $(x - 1)(2x + 4) = 0$.

Or ce trinôme du second degré a deux zéros$- 2$ et 1 et son coefficient dominant est $+ 2$, donc ce trinôme est positif sauf entre les racines $- 2$ et 1. Réponse \textbf{c.}
\item  La fonction $h$ définie sur $\R$ par $h(x)= (x + 1)\text{e}^{x}$ est :







On peut calculer $h'(x) = (x + 2)\text{e}^x$, puis $h''(x) = (x + 3)\text{e}^x$.

Comme $\text{e}^x > 0$ quel que soit le réel $x$ le signe de $h''(x)$ est celui de $x + 3$ qui est positif sur 

$]-3~;~+ \infty[$. Sur cet intervalle la fonction est convexe. Réponse \textbf{d.}

\item Une suite $\left(u_n\right)$ est minorée par 3 et converge vers un réel $\ell$.








La limite est supérieure ou égale à 3. Réponse \textbf{b.}

\item  La suite $\left(w_n\right)$ est définie par $w_1 = 2$ et pour tout entier naturel $n$ strictement positif, 

$w_{n+1} = \dfrac1n w_n$.






On peut démontrer rapidement par récurrence de :

$w_1 = 2$, \quad $w_2 = \dfrac{2}{1} = 2$, \quad $w_3 = 1, \quad w_4 = \dfrac{1}{3}w_3 = \dfrac13 = \dfrac{2}{2 \times 3} = \dfrac{2}{(4 - 1)!}$, ...

que $w_n = \dfrac{2}{(n-1)!}$, avec $n! = 1 \times 2 \times 3 \times \ldots \times n$.

Or $\displaystyle\lim_{n \to + \infty}\dfrac{2}{(n- 1)!} = 0$. Réponse \textbf{d.}
\end{enumerate}

\bigskip
